% Generated draft: requires mathematical and source review.
\documentclass[11pt]{article}
\usepackage[margin=25mm]{geometry}
\usepackage{fontspec}
\setmainfont{Latin Modern Roman}
\usepackage{amsmath,amssymb,mathtools,unicode-math}
\setmathfont{Latin Modern Math}
\usepackage{enumitem,needspace,xurl,hyperref}
\setlength{\parindent}{0pt}
\setlength{\parskip}{5pt}
\setlength{\emergencystretch}{3em}
\newcommand{\N}{\mathbb{N}}
\newcommand{\Z}{\mathbb{Z}}
\newcommand{\Q}{\mathbb{Q}}
\newcommand{\R}{\mathbb{R}}
\newcommand{\C}{\mathbb{C}}
\newcommand{\F}{\mathbb{F}}
\newcommand{\A}{\mathbb{A}}
\newcommand{\PP}{\mathbb{P}}
\newcommand{\E}{\mathbb{E}}
\newcommand{\eps}{\varepsilon}
\newcommand{\dd}{\,\mathrm{d}}
\begin{document}
\section*{UnsolvedMath ID: 20002077}
\textbf{Source problem code:} catalogue entry
\subsection{Definitions and mathematical statement}
Use \(z_j=p_j+iq_j\), \(\omega_0=\sum_j\dd p_j\wedge\dd q_j\), and the closed unit ball \(B^4=\{|z_1|^2+|z_2|^2\leq1\}\). The corrected Clifford torus is
\[
L=\{|z_1|^2=|z_2|^2=1/2\}\subset\partial B^4.
\]
The source asks whether \(L\) is weakly boundary rigid: must every Hamiltonian-isotopic torus \(K\) contained in \(B^4\) lie entirely in \(\partial B^4\)? The aid uses ambient Hamiltonian images \(K=\phi_H^1(L)\) in \(\R^4\), with only the endpoint containment required. Requiring \(\phi_H^1\) itself to be a diffeomorphism of the closed ball would preserve the boundary automatically and is not this nontrivial interpretation. The historical fact that \(K\) cannot be entirely in the open ball is weaker than the desired whole-boundary conclusion.

The aid attributes an affirmative four-dimensional answer to the peer-reviewed extremal-torus theorem of Dimitroglou Rizell. It reports a July 2026 preprint of Faisal--Li giving the analogue in every ellipsoid and dimension. Fix an action normalization
\[
E(a_1,\ldots,a_n)=\left\{z:\sum_{j=1}^n\frac{\pi|z_j|^2}{a_j}\leq1\right\},\quad a_j>0,
\quad d=\left(\sum_ja_j^{-1}\right)^{-1},
\]
\[
T_d=\{z:\pi|z_j|^2=d\text{ for every }j\}\subset\partial E.
\]
The reported weak-rigidity statement is \(\phi_H^1(T_d)\subset E\Rightarrow\phi_H^1(T_d)\subset\partial E\). Equivalently, any image not contained in \(\partial E\) has a point satisfying \(\sum_j\pi|z_j|^2/a_j>1\): strict escape, not merely failure of interior containment.

The aid further reports actual boundary rigidity when the reciprocal parameters are \(\Q\)-linearly independent:
\[
\sum_j\frac{m_j}{a_j}=0,\quad m_j\in\Z
\quad\Longrightarrow\quad m_1=\cdots=m_n=0.
\]
Under this nonresonance assumption, any Hamiltonian image of \(T_d\) contained in \(E\) equals \(T_d\) setwise. It also includes endpoints of exact Lagrangian isotopies: smooth Lagrangian embeddings \(\iota_t:T^n\to\R^{2n}\), beginning at the inclusion, with \([\iota_t^*\lambda_0]\) constant, where \(\lambda_0=\tfrac12\sum_j(p_j\dd q_j-q_j\dd p_j)\). The exactness/ambient-extension conventions must be retained; ordinary Lagrangian isotopy is not substituted.

Part (b) additionally asks to investigate other ellipsoid tori and nonremovable intersections. The balanced-product result and nonresonant refinement do not classify all tori or all such intersections. The preprint date and differing publication status of the reported results are material, and no current independent literature verification is claimed.
\subsection{Short English statement}
Must a Hamiltonian image of the Clifford torus that fits in the closed ball stay entirely on its boundary? The supplied aid reports an affirmative published four-dimensional result, a July 2026 ellipsoid preprint, and a stronger nonresonant refinement. The wider nonremovable-intersection question remains separate.
\subsection{Sources}
\begin{itemize}
\item Question 1.19(a)--(b), AIM workshop \emph{Holomorphic curves in contact geometry}: \url{https://aimath.org/WWN/contactgeom2/contactgeom2.pdf}. The source refers to Eliashberg's end-of-day discussion; the Alex Ivrii transition introduces subsequent questions.
\item The aid cites Dimitroglou Rizell's peer-reviewed four-dimensional extremal-torus theorem and a July 2026 preprint of Faisal--Li, without titles, identifiers or complete statements.
\item Local registry \texttt{data/unsolvedmath/problems.json}, permanent ID 20002077, code \texttt{catalogue entry}; canonical \texttt{aim-geometry-notes.json}, \texttt{notes[414]}; updated 14 August 2026. Its verification reports prior checking of source PDF page 6 to repair the damaged torus formula; the original literature field is empty.
\item Attempt 1: partial result, partially solved status, rigor label proved, novelty confidence low. The machine-generated aid is not peer reviewed and requires independent verification before citation.
\end{itemize}
\end{document}
